% Translated from: thesis/chapter_05_实验设计.md
% Target journal: General
% Generated: 2026-04-03

% --- Paragraph 1 ---
% # 第五章 实验设计

\chapter{Experimental Design}
\setcounter{figure}{0}
\setcounter{equation}{0}
\renewcommand{\thefigure}{\thechapter-\arabic{figure}}
\renewcommand{\theequation}{\thechapter-\arabic{equation}}

% --- Paragraph 2 ---
% ## 5.1 实验目标与验证层次

\section{Experimental Objectives and Validation Layers}

% --- Paragraph 3 ---
% 围绕第三章提出的 SDRPV 方法，本文将实验设计组织为三个相互衔接的验证层次。第一层是 baseline benchmark，用于回答 SDRPV 在统一搜索预算下相对于传统搜索基线和神经价值基线是否具有实际对局优势；第二层是 ablation benchmark，用于分析 teacher 监督与 residual 学习两个关键组成部分对最终性能的贡献；第三层是离线指标评估，用于从监督拟合角度检验模型是否真正学到了相对于基线搜索值的有效修正。

The experiments are organized into three connected validation layers around SDRPV. The first layer is the baseline benchmark, which tests whether SDRPV has practical online-match advantages over traditional search baselines and neural-value baselines under a unified search budget. The second layer is the ablation benchmark, which analyzes the contributions of teacher supervision and residual learning. The third layer is offline metric evaluation, which verifies whether the model learns effective corrections over baseline search values.

% --- Paragraph 4 ---
% ## 5.2 实验环境与数据来源

\section{Experimental Environment and Data Sources}

% --- Paragraph 5 ---
% 本文的实验环境建立在统一的 GGP 状态机之上，所有方法共享相同的规则执行机制、状态转移接口与终局评分规则。在线实验采用的游戏包括 \texttt{breakthrough}、\texttt{connectFour} 和 \texttt{hex}，它们共同覆盖了不同的棋盘规模、行动空间与博弈结构，因此能够较为全面地考察本文方法在多种棋盘型 GGP 任务中的适用性。

All experiments are built on a unified GGP state machine. All methods share identical rule execution, transition interfaces, and terminal scoring rules. Online experiments are conducted on \texttt{breakthrough}, \texttt{connectFour}, and \texttt{hex}, which cover different board sizes, action spaces, and game structures. Therefore, it is possible to comprehensively examine the applicability of this method in a variety of chessboard GGP tasks.

% --- Paragraph 6 ---
% 在线实验统一采用 fixed-time（固定时间）设置。具体而言，每一步决策的搜索时间预算设为 \texttt{0.5 s}，同时设置搜索迭代上限为 \texttt{120} 次；所有结果均以 \texttt{seed=242} 的运行留档为统计依据。若第 $m$ 步的实际搜索时间和搜索迭代数分别记为 $t_m$ 与 $k_m$，则实验预算约束满足

Online evaluation adopts a fixed-time setting. Each move has a search budget of \texttt{0.5 s} with at most \texttt{120} iterations, and all reported statistics are based on archived runs with \texttt{seed=242}. Let runtime and iterations at move $m$ be $t_m$ and $k_m$:

\begin{equation}
t_m \le 0.5,\qquad k_m \le 120.
\end{equation}

% --- Paragraph 8 ---
% 采用这一口径的原因在于，它更接近真实对局中“有限思考时间”的使用场景，同时也避免了不同随机种子和不同预算模式混合带来的解释复杂度。

This protocol better matches real-time decision conditions and avoids confounding effects from mixed seed settings or mixed budget modes.

% --- Paragraph 9 ---
% 训练数据由统一的数据构建流程生成。本文离线训练样本来自 \texttt{connectFour} 的自对弈轨迹，且由 \texttt{pure_mct} 与 \texttt{heuristic_mcts} 混合生成。数据生成阶段采用 \texttt{mixed_heuristic_pure} 策略，配置比例为 \texttt{heuristic_mcts:pure_mct = 0.8:0.2}；在本次 \texttt{seed=42} 的实际运行中，对应为 \texttt{159:41} 局（约 \texttt{79.5\%:20.5\%}）。采用该混合策略的目的，是在保留较强启发式轨迹的同时引入纯搜索轨迹，降低单一行为策略导致的分布偏置，使样本覆盖弱启发式到强启发式的更宽状态分布，从而提升价值模型对不同搜索风格对手的稳健性。随后，系统在上述自对弈过程中采集中间状态作为候选样本，并对同一状态分别执行两档搜索标注：一档为低成本预算，用于生成基线值 \texttt{b}；另一档为更高预算，用于生成 teacher 值 \texttt{q_t}。随后，系统依据终局得分生成归一化终局监督 \texttt{z}，并从当前状态提取阶段特征 \texttt{phi}（如回合进度、占位率和合法动作规模等），同时记录样本来源、游戏类型与步骤位置等 \texttt{metadata}。在此基础上，数据管线将 \texttt{b}、\texttt{q_t} 与 \texttt{z} 统一映射到 $[-1,1]$ 数值区间，并执行去重、异常值过滤与字段完整性检查，最终得到可用于训练的标准化样本

Training data are generated by a unified pipeline. Offline samples come from \texttt{connectFour} self-play trajectories produced by mixing \texttt{pure_mct} and \texttt{heuristic_mcts}. The data-generation policy is \texttt{mixed_heuristic_pure} with ratio \texttt{heuristic_mcts:pure_mct = 0.8:0.2}; under \texttt{seed=42}, this corresponds to \texttt{159:41} games (approximately \texttt{79.5\%:20.5\%}). 
The purpose of adopting this hybrid strategy is to introduce a pure search trajectory while retaining a strong heuristic trajectory, reduce the distribution bias caused by a single behavior strategy, and make the sample cover a wider state distribution from weak heuristic to strong heuristic, thereby improving the robustness of the value model to opponents with different search styles. Subsequently, the system collects intermediate states as candidate samples during the above-mentioned self-playing process, and performs two search annotations on the same state: one is a low-cost budget to generate a baseline value \texttt{b}; another file is a higher budget, which is used to generate the teacher value \texttt{q_t}. Subsequently, the system generates a normalized final supervision \texttt{z} based on the final score, and extracts the stage features \texttt{phi} ( such as round progress, occupancy rate and legal action scale, etc. ) from the current state, and records the sample source, game type and step position \texttt{metadata}. On this basis, the data pipeline maps \texttt{b}, \texttt{q_t} and \texttt{z} to the $[-1,1]$ numerical interval, and performs deduplication, outlier filtering and field integrity checking. Finally, a standardized sample that can be used for training is obtained:

\begin{equation}
x=\bigl(s,b,q_t,z,\phi,\text{metadata}\bigr).
\end{equation}






% --- Paragraph 11 ---
% 经过上述清洗与标准化处理后，最终形成 \texttt{4732} 条有效样本，并划分为训练集、验证集和测试集，规模分别为 \texttt{3785}、\texttt{473} 和 \texttt{474}。若样本总数记为 $N=4732$，则三类划分占比近似为
After cleaning and normalization, we obtain \texttt{4732} valid samples, split into train/validation/test sets of \texttt{3785}, \texttt{473}, and \texttt{474}. Let total sample size be $N=4732$:

\begin{equation}
\frac{3785}{4732}\approx 0.800,\qquad
\frac{473}{4732}\approx 0.100,\qquad
\frac{474}{4732}\approx 0.100.
\end{equation}



% --- Paragraph 13 ---
% 这种构建方式使训练目标从一开始就面向本文的最终应用场景，即在有限搜索预算下，通过学习修正量提升叶节点评估质量。
This construction method makes the training target oriented to the final application scenario of this paper from the beginning, that is, under the limited search budget, the quality of leaf node evaluation is improved by learning the correction.


% --- Paragraph 14 ---
% 【表 5-1 预留：实验对象与数据来源说明表。建议列出三类游戏、fixed-time 搜索预算、随机种子、样本划分规模以及三类实验对应的数据来源与输出内容。】

\begin{table}[htbp]
\centering
\caption{Experimental objects and data sources}
\label{tab:experiment_objects_and_data_sources}
\begin{tabular}{p{2.4cm}p{5.0cm}p{4.1cm}}
\toprule
Item & Specific setting or data source & Corresponding outputs \\
\midrule
Online games & \texttt{breakthrough}, \texttt{connectFour}, \texttt{hex} & Cross-game fixed-time match results \\
Search budget & Per move $0.5\,\mathrm{s}$, at most $120$ iterations & Unified protocol for baseline and ablation benchmarks \\
Random seeds & \texttt{seed=42} for data construction, \texttt{seed=242} for online archives & Reproducible data generation and evaluation statistics \\
Training data source & \texttt{connectFour} self-play, \texttt{mixed\_heuristic\_pure}, \texttt{heuristic\_mcts:pure\_mct=159:41} & Standardized SDRPV training samples \\
Split scale & Total $4732$, train/val/test = $3785/473/474$ & Residual training, model selection, offline evaluation \\
baseline benchmark & SDRPV vs. \texttt{pure\_mct}, \texttt{heuristic\_mcts}, \texttt{baseline\_token\_transformer} & Per-game and overall win rates \\
ablation benchmark & Full model, no-residual variant, no-teacher variant & Component-contribution analysis \\
Offline evaluation & Predictions, baseline values, and teacher values on test split & MAE, correlations, and error-gain metrics \\
\bottomrule
\end{tabular}
\end{table}

% --- Paragraph 16 ---
% ## 5.3 模型角色与对比基线设置
\section{Model Roles and Baseline Configuration}


% --- Paragraph 17 ---
% 本文在实验中考察的方法可分为两组。第一组是本文的主方法，即基于 SDRPV 的神经价值增强搜索。该方法采用棋盘 token 表示和 Transformer 价值网络，以残差学习方式估计相对于低成本搜索值的修正量，并在在线搜索中将该估计作为叶节点评估的补充信息。
Methods are divided into two groups. The first is our main method: SDRPV-based neural-value-enhanced search. It uses board-token representation and a Transformer value network, learns residual corrections over low-cost search values, and injects these corrections into leaf-node evaluation online.


% --- Paragraph 18 ---
% 第二组是对比基线，包括 \texttt{pure_mct}、\texttt{heuristic_mcts} 和 \texttt{baseline_token_transformer}。其中，\texttt{pure_mct} 代表不使用启发式信息和神经评估的纯 UCT 搜索；\texttt{heuristic_mcts} 在统一 MCTS 内核上引入历史统计驱动的动作采样与 rollout 引导，可视为更强的启发式搜索基线 \cite{finnsson2008simulation}；\texttt{baseline_token_transformer} 则代表与本文主方法具有相同表示范式、但不采用 SDRPV 残差学习目标的神经价值基线，其设置有助于辨别本文收益究竟来自统一表示本身，还是来自 SDRPV 的学习目标与搜索接入方式。
The second group contains baselines: \texttt{pure_mct}, \texttt{heuristic_mcts}, and \texttt{baseline_token_transformer}. \texttt{pure_mct} is vanilla UCT without heuristics or neural evaluation. \texttt{heuristic_mcts} adds history-guided sampling and rollout guidance \cite{finnsson2008simulation}. \texttt{baseline_token_transformer} shares the same representation style as our method but does not use SDRPV residual objectives, enabling us to isolate gains from objective design and integration strategy rather than representation alone.


% --- Paragraph 19 ---
% 这样的基线组合具有明确分工：\texttt{pure_mct} 用于考察本文方法相对于最基础搜索的提升幅度，\texttt{heuristic_mcts} 用于检验本文方法对较强启发式搜索是否仍具有竞争力，\texttt{baseline_token_transformer} 则用于评估 SDRPV 相对于已有神经价值方法的增益。由此，baseline benchmark 可以从搜索、启发式和神经价值三个角度对本文方法进行完整定位。
This baseline set has clear functional roles: \texttt{pure_mct} measures gains over minimal search, \texttt{heuristic_mcts} tests competitiveness against stronger heuristic search, and \texttt{baseline_token_transformer} measures gains over existing neural-value baselines. Therefore, the baseline benchmark can completely locate the proposed method from three perspectives: search, heuristic and neural value.


% --- Paragraph 20 ---
% ## 5.4 Baseline Benchmark 设计
\section{Baseline Benchmark Design}


% --- Paragraph 21 ---
% baseline benchmark 的目的是在统一预算和统一环境下比较 SDRPV 与三类基线的真实对局表现。实验在 \texttt{breakthrough}、\texttt{connectFour} 和 \texttt{hex} 三个游戏上进行，所有方法均采用 fixed-time 口径，即每步 \texttt{0.5 s} 搜索时间预算与 \texttt{120} 次搜索迭代上限，并以 \texttt{seed=242} 作为留档结果的统计依据。为保证分游戏结果的稳定性，每个游戏下的每组对抗配置均运行 \texttt{50} 局对弈。若记胜局数、平局数和总对局数分别为 $N_{\mathrm{win}}$、$N_{\mathrm{draw}}$ 和 $N_{\mathrm{all}}$，则本文采用的胜率统计可写为
The baseline benchmark compares SDRPV with three baselines under identical budgets and environments. Experiments run on \texttt{breakthrough}, \texttt{connectFour}, and \texttt{hex} with fixed-time settings (\texttt{0.5 s} per move, up to \texttt{120} iterations) and \texttt{seed=242} for archived statistics. Each game-method pairing is evaluated over \texttt{50} matches. Let wins, draws, and total matches be $N_{\mathrm{win}}$, $N_{\mathrm{draw}}$, and $N_{\mathrm{all}}$:

\begin{equation}
\mathrm{WinRate}
=
\frac{N_{\mathrm{win}}+0.5N_{\mathrm{draw}}}{N_{\mathrm{all}}}.
\end{equation}


% --- Paragraph 23 ---
% 在该实验中，本文方法采用 SDRPV 驱动的神经价值增强搜索；对比方法依次为 \texttt{pure_mct}、\texttt{heuristic_mcts} 和 \texttt{baseline_token_transformer}。前三者共同构成从纯搜索到启发式搜索再到神经价值搜索的递进对照，因此该 benchmark 可以清楚回答两个问题：第一，SDRPV 是否优于传统搜索基线；第二，SDRPV 相对于已有神经价值方法是否具有额外收益。
In this experiment, the proposed method uses SDRPV-driven neural value enhancement search; the comparison methods are \texttt{pure_mct}, \texttt{heuristic_mcts} and \texttt{baseline_token_transformer}. The first three together constitute a progressive control from pure search to heuristic search and then to neural value search. Therefore, this benchmark answers two questions: whether SDRPV is superior to traditional search baselines, and whether SDRPV yields additional gains over an existing neural-value baseline.


% --- Paragraph 24 ---
% 除以 SDRPV 为核心的主 baseline benchmark 外，本文还新增了一个补充 baseline benchmark，即固定 \texttt{baseline_token_transformer} 为对局主体，仅与 \texttt{pure_mct} 和 \texttt{heuristic_mcts} 进行对弈。该补充实验与主 benchmark 共用完全一致的 fixed-time 设置、随机种子与对局轮数，因此能够在“相同表示范式、相同传统搜索对手、相同预算约束”的条件下，直接比较 SDRPV 与神经价值基线对传统搜索基线的实际压制能力。换言之，主 benchmark 用于证明 SDRPV 整体更强，而补充 benchmark 则进一步说明这种优势不仅来自 token Transformer 表示带来的普遍收益，也来自 SDRPV 的残差学习目标及其搜索接入方式。
We also add a supplemental baseline benchmark by fixing \texttt{baseline_token_transformer} as the main agent and playing only against \texttt{pure_mct} and \texttt{heuristic_mcts} under exactly the same fixed-time settings, seeds, and game counts. This allows a direct comparison of suppression strength against traditional-search opponents under matched conditions. In other words, the main benchmark is used to prove that SDRPV is stronger as a whole, while the supplementary benchmark further shows that this advantage not only comes from the general benefits brought by token Transformer representation, but also from the residual learning goal of SDRPV and its search access method.


% --- Paragraph 25 ---
% 【图 5-1 预留：baseline benchmark 组织流程图。建议展示“模型加载 -> 搜索代理配置 -> fixed-time 对局执行 -> 单局结果汇总 -> 分游戏统计与总体统计”的流程。】
\begin{figure}[htbp]
\centering
\includegraphics[width=0.90\linewidth]{figures/chapter5/baseline_benchmark_pipeline.pdf}
\caption{Pipeline of the baseline benchmark.}
\label{fig:baseline_benchmark_pipeline}
\end{figure}

% --- Paragraph 27 ---
% 【表 5-2 预留：主 baseline benchmark 与补充 baseline benchmark 的关系表。建议列出两组实验的 agent A、对手集合、控制变量以及各自回答的研究问题。】

\begin{table}[htbp]
\centering
\caption{Design relation between the main and supplemental baseline benchmarks}
\label{tab:baseline_and_supplement_design}
\begin{tabular}{p{2.8cm}p{3.2cm}p{4.0cm}p{4.2cm}}
\toprule
Experiment group & Agent A & Opponent set & Main question addressed \\
\midrule
Main baseline benchmark & SDRPV & pure\_mct, heuristic\_mcts, baseline\_token\_transformer & Whether SDRPV is superior overall under unified budgets \\
Supplemental baseline benchmark & baseline\_token\_transformer & pure\_mct, heuristic\_mcts & Whether SDRPV gains exceed those from neural baseline representation alone \\
\bottomrule
\end{tabular}
\end{table}
% --- Paragraph 29 ---
% 【图 5-2 预留：主 baseline benchmark 与补充 baseline benchmark 的关系示意图。建议用双路径对照展示 SDRPV 与 \texttt{baseline_token_transformer} 分别对接 \texttt{pure_mct}、\texttt{heuristic_mcts}，并在中间标出 \texttt{SDRPV vs baseline_token_transformer} 的直接比较关系。】
\begin{figure}[htbp]
\centering
\includegraphics[width=0.88\linewidth]{figures/chapter5/baseline_supplement_relation.pdf}
\caption{Relation between the main and supplemental baseline benchmarks.}
\label{fig:baseline_supplement_relation}
\end{figure}


% --- Paragraph 31 ---
% ## 5.5 Ablation Benchmark 设计
\section{Ablation Benchmark Design}

% --- Paragraph 32 ---
% ablation benchmark 的目的是拆解 SDRPV 中两个关键设计的作用，即较强搜索监督信号与残差学习目标。为此，实验比较完整方案、去除残差项的方案以及去除 teacher 监督的方案，借此分析最终性能提升究竟来自哪一部分设计，以及两者之间是否存在协同效应。
The ablation benchmark isolates the effects of two key designs in SDRPV: strong-search teacher supervision and residual learning objectives. We compare the full model, a no-residual variant, and a no-teacher variant to identify each component's contribution and possible synergy.


% --- Paragraph 33 ---
% 该实验与 baseline benchmark 保持完全一致的环境控制方式，同样在 \texttt{breakthrough}、\texttt{connectFour} 和 \texttt{hex} 上运行，同样采用每步 \texttt{0.5 s} 的 fixed-time 搜索预算与 \texttt{120} 次迭代上限，并同样以 \texttt{seed=242} 的结果作为统计依据。每个游戏下的每组配置均运行 \texttt{50} 局，以保证消融比较具有足够的稳定性。
The benchmark uses the same control settings as the baseline benchmark to ensure that the ablation has sufficient stability: \texttt{breakthrough}, \texttt{connectFour}, and \texttt{hex}; fixed-time budget of \texttt{0.5 s} and \texttt{120} iterations; \texttt{seed=242}; and \texttt{50} matches per configuration in each game.

% --- Paragraph 34 ---
% 通过这一设置，ablation benchmark 可以回答如下问题：如果不学习残差修正量，而是直接依赖较强搜索值进行监督，在线表现会发生怎样的变化；如果保留残差框架但移除 teacher 监督信号，模型是否仍能形成有效的搜索增强。也就是说，消融实验是用来解释 SDRPV 的性能来源和结构合理性。
Through this setting, the ablation benchmark can answer the following questions: if we do not learn the residual correction, but directly rely on the strong search value for supervision, how will the online performance change; if the residual framework is retained but the teacher supervision signal is removed, whether the model can still form an effective search enhancement. That is to say, the ablation experiment is used to explain the performance source and structural rationality of SDRPV.

% --- Paragraph 35 ---
% 【表 5-3 预留：ablation 配置说明表。建议列出完整方案、去除残差方案、去除 teacher 方案三类配置在监督目标、在线接入方式和预期作用上的区别。】
\begin{table}[htbp]
\centering
\caption{Ablation configurations}
\label{tab:ablation_configurations}
\begin{tabular}{p{3.2cm}p{3.5cm}p{3.6cm}p{3.2cm}}
\toprule
Configuration & Supervision objective & Online integration style & Expected role \\
\midrule
\texttt{full\_residual\_teacher} & Learn teacher-guided residual correction over \texttt{b} & Residual value in value-full search & Full mainline model with both teacher and residual design \\
\texttt{ablation\_no\_residual\_teacher} & Directly fit teacher values without residuals & Keep neural integration but remove \texttt{baseline + correction} form & Test contribution of residual design \\
\texttt{ablation\_no\_teacher\_residual} & Learn residual objective without teacher distillation & Keep residual integration but remove strong-search supervision & Test contribution of teacher supervision direction \\
\bottomrule
\end{tabular}
\end{table}


% --- Paragraph 37 ---
% ## 5.6 离线指标与控制变量说明
\section{Offline Metrics and Controlled Variables}


% --- Paragraph 38 ---
% 离线指标实验的作用，是从监督学习角度验证残差模型是否真正学到了比低成本基线搜索值更接近较强搜索值的估计。本文主要报告相关性、排序相关性、相对于较强搜索值的绝对误差，以及相对于基线搜索值的误差改进幅度等指标。若测试集上的预测值记为 $\hat{v}_i$，对应 teacher 值记为 $q_{t,i}$，基线搜索值记为 $b_i$，则核心指标可写为
Offline evaluation verifies whether residual models produce estimates closer to strong-search targets than baseline values. Let prediction be $\hat{v}_i$, teacher target be $q_{t,i}$, and baseline value be $b_i$ on test samples. Core metrics are



\begin{equation}
\mathrm{MAE}_{\hat{v}}
=
\frac{1}{N}\sum_{i=1}^{N}\lvert \hat{v}_i-q_{t,i}\rvert,
\end{equation}



\begin{equation}
\mathrm{MAE}_{b}
=
\frac{1}{N}\sum_{i=1}^{N}\lvert b_i-q_{t,i}\rvert,
\end{equation}



\begin{equation}
\mathrm{Gain}_{\mathrm{MAE}}
=
\mathrm{MAE}_{b}-\mathrm{MAE}_{\hat{v}},
\end{equation}

% --- Paragraph 42 ---
% 以及
and correlation metrics



\begin{equation}
\rho_{\mathrm{P}}
=
\mathrm{corr}(\hat{v},q_t),\qquad
\rho_{\mathrm{S}}
=
\mathrm{corr}\bigl(\mathrm{rank}(\hat{v}),\mathrm{rank}(q_t)\bigr).
\end{equation}


% --- Paragraph 44 ---
% 与仅报告训练损失相比，这组指标更能直接对应第三章提出的问题定义，即模型是否学会了有效修正量而不是简单拟合某个静态标签。
Compared with only reporting training loss, this set of indicators can more directly correspond to the definition of the problem proposed in Chapter 3, that is, whether the model has learned effective corrections rather than simply fitting a static label.

% --- Paragraph 45 ---
% 为了使 baseline benchmark、ablation benchmark 与离线指标分析之间具备可比性，本文对关键变量进行了统一控制。首先，所有在线实验共享同一状态机、同一 MCTS 主循环和相同的终局评分规则，不同方法之间仅改变叶节点评估策略。其次，所有在线实验统一采用 \texttt{breakthrough}、\texttt{connectFour}、\texttt{hex} 三个游戏，统一采用 fixed-time 口径，即每步 \texttt{0.5 s} 搜索预算与 \texttt{120} 次搜索迭代上限，统一采用 \texttt{seed=242} 和每组配置 \texttt{50} 局的设置。再次，神经价值基线与本文方法都在相同的搜索框架下接入 MCTS，从而保证比较对象是“价值评估质量及其搜索增益”，而不是外围实现差异。
To ensure comparability across baseline benchmark, ablation benchmark, and offline evaluation, key variables are controlled uniformly: same state machine, same MCTS loop, same terminal scoring; same games (\texttt{breakthrough}, \texttt{connectFour}, \texttt{hex}); same fixed-time settings (\texttt{0.5 s}, \texttt{120} iterations); same \texttt{seed=242} and \texttt{50} matches per configuration; and matched MCTS integration styles for neural baselines and SDRPV.









